% Orrery: EMNLP System Demonstration paper draft
% Target: ACL/EMNLP demo track (6 pages + references, ACL rolling style)
% Compile with the acl.sty / ACL LaTeX template (https://github.com/acl-org/acl-style-files)
%
% CONVENTIONS FOR CO-AUTHORS
%   % TODO markers flag numbers/claims that are NOT yet backed by committed
%   code or saved results. Do not leave any % TODO in the submitted PDF.
%   Verified figures (from documentation/PROJECTION_FIDELITY.md and the code)
%   are used as-is. See memory: project_orrery_paper_claims for the audit.

\documentclass[11pt]{article}
\usepackage[]{ACL2023}   % <-- replace with the current cycle's style file
\usepackage{times}
\usepackage{latexsym}
\usepackage[T1]{fontenc}
\usepackage[utf8]{inputenc}
\usepackage{microtype}
\usepackage{graphicx}
\usepackage{booktabs}
\usepackage{amsmath}
\usepackage{url}

\title{Orrery: An Interactive Platform Unifying Embedding Visualisation,
Topic Extraction, and Sparse-Autoencoder Interpretability}

\author{
  Anonymous \\
  Affiliation \\
  \texttt{email@domain} \\
}

\begin{document}
\maketitle

%======================================================================
\begin{abstract}
Orrery is an open-source platform for interactive visualisation of
embedding spaces with native support for Sparse Autoencoder (SAE)
interpretability. Given any local or HuggingFace dataset, it produces an
explorable 3D constellation through a single pipeline of embedding,
projection, clustering, automatic topic labelling, and---optionally---the
collection of per-document SAE feature activations. The resulting map can
be searched visually, semantically, lexically (BM25), or by SAE feature
name. SAE decoder vectors are visualised as a navigable constellation of
their own, where entering a document highlights the features it activates.
Individual features can be inspected in a Neuronpedia-style dashboard and
applied directly to steer the model. The interface sustains 60\,fps on
250k points---including animations and nebula shaders---on 8\,GB of RAM.
Orrery unifies corpus-scale semantic exploration with internal feature
interpretation and causal model intervention, and is released under the
Apache~2.0 license at \url{[URL]}.
\end{abstract}

%======================================================================
\section{Introduction}
\label{sec:intro}

Researchers working with text embeddings face fragmented tooling.
Interactive visualisation (TensorBoard Embedding Projector
\citep{smilkov2016projector}), topic analysis (BERTopic
\citep{grootendorst2022bertopic}), and mechanistic interpretability
(Neuronpedia \citep{neuronpedia}) exist as separate systems. Moving
between them requires data export, format conversion, and constant context
switching. No existing tool lets a researcher go from ``\emph{what does my
embedding space look like?}'' to ``\emph{which model features drive this
structure?}'' within a single workflow.

\paragraph{Contribution.} We present \textbf{Orrery}, an open-source
platform that integrates three capabilities that are usually siloed:
(i)~interactive 2D/3D embedding visualisation that sustains 250k+ points at
interactive frame-rates on commodity hardware; (ii)~an automatic topic
extraction pipeline with hierarchical reduction and LLM labelling; and
(iii)~Sparse-Autoencoder analysis with live model inference, per-document
feature activation, and interactive steering. To our knowledge Orrery is
the first system to connect general-purpose embedding visualisation to
mechanistic-interpretability SAE analysis, closing the loop:
\emph{find a semantic pattern in the map $\rightarrow$ identify candidate
features or directions $\rightarrow$ inspect the evidence $\rightarrow$
intervene in the model $\rightarrow$ observe the behavioural effect.}

\paragraph{Emergent structure as validation.} To demonstrate that the
platform surfaces genuine scientific insight, we report two findings that
its analytical-colouring workflow makes directly visible. First, the
embedding of colour \emph{names} self-organises around perceptual colour
space: a UMAP-3D layout of the XKCD colour survey is \emph{more}
colour-coherent (Mantel $\rho = 0.60$) than the 3072-d embedding it was
projected from ($\rho = 0.29$; \S\ref{sec:xkcd}). Second, psycholinguistic
dimensions such as concreteness and imageability are linearly decodable
from the embedding geometry, appearing as visible spatial gradients
(\S\ref{sec:concreteness}). Finally, we validate the SAE machinery by
recovering an ActAdd-style refusal direction on Gemma-3-4b-it
(\S\ref{sec:refusal}).

%======================================================================
\section{Related Work}
\label{sec:related}

\paragraph{Embedding visualisers.} The TensorBoard Embedding Projector
\citep{smilkov2016projector} pioneered interactive PCA/t-SNE/UMAP browsing
but consumes only precomputed vectors and caps at tens of thousands of
points. Embedding Atlas \citep{embeddingatlas}, Nomic Atlas
\citep{nomicatlas}, and Latent Scope \citep{latentscope} scale further and
add clustering, but Atlas and Nomic are closed or cloud-hosted and none
connect the map to a model's internal features. Artistic work such as
Google's \emph{Embedding Galaxy}\footnote{\url{https://artsexperiments.withgoogle.com/tsnemap/}}
inspires our ``star cartography'' rendering but is not a research tool.

\paragraph{Topic modelling.} Orrery's topic pipeline follows BERTopic
\citep{grootendorst2022bertopic}: HDBSCAN over a UMAP reduction, followed
by class-based TF--IDF (c-TF--IDF). We reuse BERTopic's LLM labelling
prompt directly. Our contribution here is integration, not algorithm.

\paragraph{SAE interpretability.} Sparse Autoencoders decompose model
activations into interpretable features \citep{bricken2023monosemanticity,
templeton2024scaling}. Neuronpedia \citep{neuronpedia} hosts feature
dashboards, and SAEDashboard/SAE-vis provide static visualisations. We use
the Gemma Scope SAEs \citep{lieberum2024gemmascope} and the Neuronpedia
label/activation data, but---unlike prior tools---place features in the
same interactive space as documents and expose live steering. Our steering
follows ActAdd \citep{turner2023actadd}; the refusal experiment adapts
\citet{arditi2024refusal}.

%======================================================================
\section{System Design}
\label{sec:design}

\paragraph{Datasets vs.\ collections.} Orrery separates two concepts. A
\emph{dataset} is a set of documents (text, image, audio) stored and
indexed in DuckDB. Each dataset can have one or more \emph{collections}: a
collection is a set of vectors (in ChromaDB), a set of projections (in
DuckDB), and optionally a set of topics and a set of per-document SAE
activations (both in DuckDB). Decoupling the two lets a user embed the same
dataset many times---different prompts, column combinations, or
models---without duplicating the documents, which was the main limitation
of a vector-store-only design.

\paragraph{Dual-database architecture.} DuckDB
(\texttt{main.duckdb}) is the central orchestrator: documents, JSON
metadata, native \texttt{FLOAT[]} projection arrays, normalised topic
tables, and sparse SAE activation rows. ChromaDB stores only IDs plus dense
vectors, serving approximate-nearest-neighbour similarity search (cosine).
DuckDB handles the relational and columnar-scan workloads (metadata
filters, \texttt{ILIKE} text search, FTS/BM25); ChromaDB handles ANN. A
\texttt{vector\_collections} table links one dataset to its many
embeddings. Figure~\ref{fig:arch} shows the data flow.

\begin{figure}[t]
  \centering
  % TODO Fig 2: architecture diagram
  % Data Sources -> Embedding Providers -> DuckDB + ChromaDB
  %              -> Topic Extraction -> GraphQL API -> Frontend
  \fbox{\parbox{0.9\columnwidth}{\centering\vspace{2em}
  \textit{[architecture diagram placeholder]}\vspace{2em}}}
  \caption{Data flow: sources $\rightarrow$ embedding providers
  $\rightarrow$ dual storage (DuckDB + ChromaDB) $\rightarrow$ topic
  extraction $\rightarrow$ GraphQL API $\rightarrow$ WebGL frontend.}
  \label{fig:arch}
\end{figure}

\subsection{The embedding pipeline}
\label{sec:pipeline}
Starting from a dataset loaded from HuggingFace or a local file (JSON,
Parquet, CSV) and shown as a preview, the workflow has one required step
(projections) and three optional ones.

\begin{enumerate}[nosep]
  \item \textbf{Embed.} The user chooses a column combination and a
  template, referencing columns by name in braces and keeping the prompt
  outside them, e.g.\ \texttt{\{word\}: \{definition\}}. Embedding is
  computed by one of seven providers---SentenceTransformers (local
  default), Ollama, BGE/FlagEmbedding, Qwen (all local), or the OpenAI,
  Cohere, HuggingFace, and Gemini APIs---with task-type support for Gemma
  and Gemini and hardware auto-detection (MPS/CUDA/CPU). Vectors go to
  ChromaDB; selected metadata columns go to DuckDB. This step can be
  \emph{skipped} for datasets that already carry a vector column.
  \item \textbf{Project (required).} We automatically compute and store 2D
  and 3D projections for both PCA and UMAP.
  \item \textbf{Topic model (optional).} \S\ref{sec:topics}.
  \item \textbf{SAE activations (optional).} \S\ref{sec:sae}.
\end{enumerate}

\subsection{Topic extraction}
\label{sec:topics}
When enabled, a BERTopic-style pipeline runs. Because BERTopic is modular
we expose several clustering algorithms---HDBSCAN (default), K-means, GMM,
and Spectral clustering (suited to small datasets, $O(n^2)$
memory)---applied by default to a fresh 5-component UMAP of the raw
vectors. Labels come from c-TF--IDF keywords, optionally upgraded with LLM
labels (Gemini/OpenAI) using BERTopic's prompt, with incremental saves and
resume. An optional second reduction step merges topics to a target count
while preserving the originals as subtopics, enabling hierarchical
Tableau-style colouring.
% TODO: cluster-quality metrics (silhouette / Davies-Bouldin) are described
% in the outline but a standalone evaluator (DBCV, silhouette, C_v/U_Mass)
% lives in interpretability_backend/evaluation and is NOT yet persisted /
% exposed. Either wire it in and report numbers, or omit the claim.

\subsection{Colour support}
\label{sec:colour}
Any metadata column can drive colouring. Orrery supports categorical,
sequential, diverging, and monochrome scales, including Crameri's
scientific colour maps \citep{crameri2020colourmaps} for perceptually
uniform, colour-vision-deficiency-safe rendering---important both for
interpretability research and for displaying historical corpora. When a
dataset carries a hex colour column, embedding auto-detects it and derives
a numeric \texttt{mapped\_colour} field via a Hilbert-curve strip that
reduces a 3D colour space to a single perceptually-ordered dimension
(\S\ref{sec:xkcd}, Appendix). Colour schemes round-trip through the URL and
can be saved per collection as a default.

%======================================================================
\section{Visualisation and Interaction}
\label{sec:viz}

\paragraph{Rendering.} The scatter plots are WebGL-accelerated Plotly.js
(2D and 3D). We fork Plotly.js to fix an $O(n\cdot m)$ trace-update cost
that made real-time overlays unusable on large collections; with the fix,
250k+ points render at 60\,fps on 8\,GB of RAM, with marker size and
opacity scaled logarithmically to point count.

\paragraph{Search as filter.} A collection can be searched four ways, with
results highlighted as a constellation whose stars grow \emph{warmer} with
similarity: (i)~\textbf{semantic} search (also triggered by clicking any
point, which flies the camera in and retrieves top-$k$ neighbours in the
original space); (ii)~\textbf{text} search (server-side DuckDB
\texttt{ILIKE}/BM25 with per-column, exact/partial, and case controls);
(iii)~\textbf{SAE feature-name} search (\S\ref{sec:sae}); and
(iv)~\textbf{prompt} search on SAE collections.

\paragraph{Filtering.} The legend panel doubles as a filter: a categorical
colour field renders as a searchable, scrollable, counted list whose
entries toggle points on and off; a numerical field renders as a
colour-scaled histogram with draggable handles that set the scale, and the
analytics panel offers a second histogram whose handles graphically filter
a range. An advanced panel composes explicit predicates (equals, not-equals,
in, not-in, range). Temporal fields are auto-detected and filtered through a
draggable range picker.

\paragraph{Rendering features for cluster identification.} In 3D, an
optional \emph{nebula} mode draws translucent halos around topic clusters
via camera-synced sprite shaders on a blended overlay canvas, and a
collision-avoidance system projects 3D label anchors to screen space (using
MVP matrices read from Plotly's WebGL internals) and greedily packs
non-overlapping labels by a similarity-weighted priority. A deferred
selected-point pattern keeps camera fly-to animations off the $O(n)$ Plotly
diff path. A one-click screenshot exports the composited scene (plot, haze,
labels, legend) at device resolution.

%======================================================================
\section{Sparse-Autoencoder Support}
\label{sec:sae}

SAE support is the platform's primary novel contribution, and it comprises
two workflows: visualising a \emph{space of features}, and collecting
features \emph{per document}.

\paragraph{From-scratch SAE core.} Rather than wrapping SAELens or
TransformerLens, Orrery implements JumpReLU and TopK SAEs directly
(\raise.17ex\hbox{$\scriptstyle\sim$}135 lines). A hook manager attaches
SAEs to specific sites (\texttt{resid\_post}, \texttt{attn\_out},
\texttt{mlp\_out}) and can co-attach multiple SAEs---including two widths at
one layer---in a single forward pass, with prefill-only semantics and
steering composition. Mid-layer activation capture uses a forked
\texttt{gemma\_pytorch} internal cache rather than standard forward hooks.
An in-memory SAE weight cache avoids re-reading the
\raise.17ex\hbox{$\scriptstyle\sim$}320\,MB \texttt{params.safetensors} on
repeated loads, improving chat time-to-first-token.
% Honest scope: the from-scratch SAE machinery is model-agnostic, but the
% live InterpretService inference wrapper is currently Gemma-only.

\subsection{Visualising SAE collections}
The feature space of an SAE can be laid out two ways. First, each feature's
decoder row (dimension $R$, the residual-stream width) is treated as a
vector and projected. Second, each feature is labelled by an
autointerpreter and the label is embedded with a conventional encoder.
Either route yields a \{\texttt{feature\_index, label, vector, top\_logit,
documents}\} record that projects into 3D like any collection. For any
Gemma Scope SAE, Orrery can download decoder vectors and metadata from the
Neuronpedia S3 bucket and ingest them into DuckDB/ChromaDB in one click.

A projected SAE collection gains two abilities over an ordinary one.
(i)~\textbf{Prompt highlighting:} entering a prompt runs it through the
model, max-pools each feature's activation across tokens, and highlights
and ranks the most-activated features directly on the scatter plot.
(ii)~\textbf{Cross-linking:} right-clicking a feature (found by vector or
label search) opens it in the Feature Explorer for inspection and steering.

\subsection{Per-document SAE activations}
For an ordinary collection, the user can additionally collect SAE
activations: every document is run through Gemma with SAE hooks, and for
each feature we keep the max activation across tokens, stored as sparse
\((\texttt{item\_id, feature\_index, activation})\) rows
(\raise.17ex\hbox{$\scriptstyle\sim$}100--150 nonzero per document). This
backs a two-hop \emph{feature-label} search: a query such as
``poetry'' \texttt{ILIKE}-matches SAE features, and documents are ranked by
their max activation over the matching features---searching a corpus by the
model's internal concepts rather than its surface words.

\subsection{Feature Explorer and steering}
A Neuronpedia-style page (token-strip activation heatmaps, top/bottom logit
bar charts) supports three query modes: text search and semantic search
over labels, and prompt search that shows, per token, which feature fires
(or the max/mean activation list over the prompt). Unlike Neuronpedia, we
keep inspection and \emph{steering} on the same page: an interesting
feature can be injected additively onto its decoder direction at a chosen
strength and the steered generation streamed token-by-token over a
WebSocket, side-by-side with a seeded baseline for comparison.
% Honest scope: steering.py implements additive/orthogonal/ablation/
% projection-cap modes, but only ADDITIVE is reachable from the demo UI.

%======================================================================
\section{Case Studies}
\label{sec:cases}

\subsection{Colour words mirror perceptual colour space}
\label{sec:xkcd}
Nonlinear structure in the embedding space can still be read off visually.
We embed the XKCD colour survey (\raise.17ex\hbox{$\scriptstyle\sim$}950
English colour terms, each with a hex code aggregated from 200k
respondents \citep{xkcd}) with Gemini and colour each point by its true
hex value. The layout intuitively recovers the colour wheel: greens,
blues, and warm tones each pool together (Fig.~\ref{fig:xkcd}).

We quantify this with a Mantel test (rank correlation between two pairwise
distance structures) comparing the projection to a \emph{perceptual}
reference---CIEDE2000 $\Delta E$ over CIELAB---and to the original
embedding geometry (cosine). The UMAP-3D layout attains global Mantel
$\rho = 0.60$ against perceptual colour, \emph{higher} than the raw
3072-d Gemini embedding it was built from ($\rho = 0.29$), and higher than
PCA-3D ($\rho = 0.47$); all correlations are strongly significant
(permutation $z = 29$--$121$, $p_{\text{emp}} < 0.001$). Strikingly, UMAP
\emph{concentrates} a colour signal that is diffuse in the full embedding:
the projection makes visible a relationship that is hidden in the original
space. (PCA-3D, by contrast, more faithfully preserves the embedding's
\emph{global} geometry, $\rho = 0.49$ vs.\ $0.29$---the honest picture of
the vectors; UMAP wins on local neighbourhood structure.)

\begin{figure}[t]
  \centering
  % TODO Fig 3: gallery/XKCD_embedded_colourspace.png (Gemini, 3D, hex-coloured)
  \fbox{\parbox{0.9\columnwidth}{\centering\vspace{3em}
  \textit{[XKCD colour-word constellation]}\vspace{3em}}}
  \caption{XKCD colour survey embedded with Gemini, projected to UMAP-3D,
  each point coloured by its true hex value. The semantic embedding of
  colour \emph{names} self-organises around perceptual colour space.}
  \label{fig:xkcd}
\end{figure}

\subsection{Psycholinguistic dimensions as spatial gradients}
\label{sec:concreteness}
A central use of an interpretability visualiser is spotting linear
directions in the embedding space---the setting of ``mass-mean'' probing in
Geometry of Truth \citep{marks2023geometry}. We embed English words from
the Brysbaert concreteness ratings \citep{brysbaert2014concreteness} and
colour each point by its concreteness score; the abstract--concrete axis is
visibly linearly separable. We replicate this on the Glasgow Norms
\citep{scott2019glasgow}, which supply seven psycholinguistic dimensions
(arousal, valence, dominance, concreteness, imageability, familiarity,
gender association), training linear and non-linear probes on each and
testing on MiniLM and Gemma-embedding-300m.
% TODO: report actual probe R^2 per dimension/model (Glasgow CSV in
% resources/uploads is 4,682 rows, NOT 19,537; runnable). Prior draft's
% "R^2 up to 0.9" and Brysbaert "93%" are NOT yet backed by saved results.
% Fill Table~\ref{tab:probes} and the numbers below before submission.
Of these, concreteness, imageability, and valence are the most cleanly
decodable [TODO $R^2$], suggesting the embedding space self-organises along
psycholinguistic axes without supervision. The number of such linear
directions is unknown; we hope the visualiser accelerates their discovery.

\begin{table}[t]
  \centering
  \small
  % TODO Table 1: fill with real linear-probe R^2 values
  \begin{tabular}{lcc}
    \toprule
    Dimension & MiniLM & Gemma-300m \\
    \midrule
    Concreteness  & [TODO] & [TODO] \\
    Imageability  & [TODO] & [TODO] \\
    Valence       & [TODO] & [TODO] \\
    Arousal       & [TODO] & [TODO] \\
    Dominance     & [TODO] & [TODO] \\
    Familiarity   & [TODO] & [TODO] \\
    Gender assoc. & [TODO] & [TODO] \\
    \bottomrule
  \end{tabular}
  \caption{Linear-probe $R^2$ recovering Glasgow Norms dimensions from
  embedding vectors. \emph{(placeholder---fill from runs)}}
  \label{tab:probes}
\end{table}

\subsection{Digital-humanities corpora at scale}
\label{sec:dh}
To show scale and humanities use, we embed two corpora at the sentence
level: the full Lacan corpus (\raise.17ex\hbox{$\scriptstyle\sim$}150k
sentences) and a section of a Sanskrit travelogue corpus. The layout
separates texts by time period, and even at 150k+ documents with nebula
mode active the scatter plot stays fully interactive, letting scholars
search and browse the corpus visually.
% Corpus sizes standardised per memory: Lacan = 150k (committed subset is
% 41k); WordNet = 212k senses if used as an additional scale demo.

\subsection{SAE validation: refusal direction}
\label{sec:refusal}
To validate steering, we adapt \citet{arditi2024refusal} to Gemma-3-4b-it
via our \texttt{gemma\_pytorch} fork. Unlike the original---which ablates
the refusal direction geometrically at every layer---we use ActAdd
\citep{turner2023actadd}, matching how we steer SAE features. We grid-search
the layer and activation range that flips the most safety-benchmark prompts
from refused to accepted while preserving coherence.
% TODO: report the real layer and bypass rate. Draft said "layer 11 / 100%
% on HarmBench"; collection metadata says layer 14 and the default eval set
% is JailbreakBench, and the eval directory was not saved. Re-run, save the
% eval, and report the true numbers + benchmark name before asserting.
We release the experiment code and the extracted refusal vector in the
repository to support LLM-safety research.

%======================================================================
\section{Comparison with Existing Systems}
\label{sec:comparison}

\begin{table*}[t]
  \centering
  \small
  \begin{tabular}{lccccc}
    \toprule
    Feature & Orrery & TB Projector & BERTopic & Neuronpedia & Nomic Atlas \\
    \midrule
    Interactive 2D/3D viz            & \checkmark & \checkmark & --- & --- & \checkmark \\
    Custom multi-provider embedding  & \checkmark & --- (precomputed) & --- & --- & partial \\
    Topic extraction + LLM labels    & \checkmark & --- & \checkmark & --- & --- \\
    Hierarchical topic reduction     & \checkmark & --- & \checkmark & --- & --- \\
    SAE feature analysis             & \checkmark & --- & --- & \checkmark & --- \\
    SAE steering + chat              & \checkmark & --- & --- & --- & --- \\
    Scatter $\leftrightarrow$ SAE linking & \checkmark & --- & --- & --- & --- \\
    Metadata analytical colouring    & \checkmark & limited & --- & --- & limited \\
    Temporal filtering               & \checkmark & --- & --- & --- & --- \\
    Open source                      & \checkmark & \checkmark & \checkmark & partial & --- \\
    Local / offline                  & \checkmark & \checkmark & \checkmark & --- & --- \\
    Scale (points)                   & 250k+ & \raise.17ex\hbox{$\scriptstyle\sim$}50k & N/A & N/A & 1M+ (cloud) \\
    \bottomrule
  \end{tabular}
  \caption{Orrery occupies a unique intersection: the only system pairing
  general-purpose embedding visualisation with SAE mechanistic
  interpretability.}
  \label{tab:comparison}
\end{table*}

Table~\ref{tab:comparison} positions Orrery. TensorBoard Projector
visualises but cannot analyse model internals; BERTopic extracts topics but
has no interactive layer; Neuronpedia explores SAE features but cannot tie
them to embedding-level structure; Nomic Atlas scales further but is
closed-source and lacks SAE integration. Orrery is the only system that
combines all three.

%======================================================================
\section{Demonstration}
\label{sec:demo}

A 3--5 minute live walkthrough: (1)~open a pre-loaded dataset and show topic
clusters with LLM labels; (2)~recolour by different metadata fields;
(3)~run semantic search with glow-highlighted, topic-scoped results;
(4)~switch to XKCD colours and show colour self-organisation; (5)~open the
Feature Explorer and browse SAE activation heatmaps; (6)~enter a prompt and
watch per-token features light up on the scatter plot; (7)~activate a
steering preset and show the behavioural change in chat against a seeded
baseline; (8)~switch to the refusal vector to demonstrate safety-relevant
steering. A hosted instance with pre-loaded datasets (XKCD Colours, Glasgow
Norms, a WordNet subset, and a safety benchmark), a \texttt{docker compose
up} path for local reproduction, and a recorded video are available at
\url{[URL]}.

%======================================================================
\section{Limitations}
\label{sec:limitations}

Orrery's SAE inference wrapper is currently Gemma-specific, though the SAE
core is model-agnostic; extending live inference to other model families
(e.g.\ Qwen) is future work. Topic extraction follows BERTopic's approach,
so the contribution there is integration rather than a new algorithm. As
reported in the SAE literature, most individual features do not yield
cleanly interpretable steering effects. Rendering scales to
\raise.17ex\hbox{$\scriptstyle\sim$}250--500k points, below cloud systems
that reach 1M+. Only additive steering is exposed in the demo UI, and
cluster-quality metrics are computed offline rather than surfaced in the
interface.

%======================================================================
\section{Conclusion}
\label{sec:conclusion}

Orrery unifies embedding visualisation, topic extraction, and SAE
interpretability in a single open-source platform. Its analytical-colouring
workflow makes emergent structure directly visible---colour-word embeddings
align with perceptual colour space (UMAP Mantel $\rho = 0.60$, above the
raw embedding's $0.29$) and psycholinguistic dimensions surface as linear
gradients---while the SAE integration lets researchers trace that structure
back to individual model features and interactively steer behaviour, a
loop not available in any existing tool.

%======================================================================
% \section*{Ethics Statement}
% TODO: note the refusal / jailbreak steering is released for defensive
% safety research; discuss dual-use and responsible-release stance.

% \section*{Acknowledgements}

\bibliographystyle{acl_natbib}
\bibliography{orrery}   % see orrery.bib (create alongside)

\appendix
\section{Hilbert-curve colour strips}
% TODO: describe generate_color_strips.py — mapping a 3D colour space to a
% single perceptually-ordered dimension via a space-filling Hilbert curve.

\section{Additional probe results}
% TODO: full Glasgow Norms / Brysbaert probe tables, per model.

\end{document}
